\section{Protocolos de Comunicación: Gateway A2A y JSON-RPC 2.0}
\label{sec:protocols}

\subsection{Especificación del Estándar A2A (Agent-to-Agent)}
El estándar A2A establece un marco formal e interoperable para la comunicación entre sistemas multiagente autónomos. Sobre la capa de transporte seguro (HTTPS / mTLS), el intercambio de mensajes se articula mediante el protocolo de invocación remota JSON-RPC 2.0.

\subsection{Estructura de la Agent Card (/.well-known/agent.json)}
El punto de anclaje semántico del sistema es su Agent Card, expuesta públicamente en la ruta estandarizada \texttt{/.well-known/agent.json}. Esta especificación describe las habilidades (\textit{skills}) ofrecidas por el orquestador, los esquemas de entrada y salida, así como los requerimientos criptográficos de autenticación para cada operación soportada.

\subsection{Protocolo de Handshake y Negociación de Agent Cards}
Antes de dar inicio formal a una negociación comercial, los agentes y sus respectivas entidades (\textit{tenants}) deben establecer un canal seguro de confianza mutua mediante el protocolo de \textit{handshake} (apretón de manos). Este procedimiento garantiza que únicamente contrapartes autorizadas y verificadas puedan acceder al catálogo de servicios y comenzar el intercambio de propuestas.

Tal como se detalla en la Figura~\ref{fig:handshake_flow}, el flujo opera a través de las siguientes fases secuenciales:
\begin{enumerate}
    \item \textbf{Solicitud de Conexión y Validación de Identidad:} El \textit{Tenant A} manifiesta su intención de negociar con el \textit{Tenant B} a través del software orquestador. El sistema procede inmediatamente a validar las credenciales de identidad del \textit{Agente A}. En caso de discrepancia criptográfica o firma inválida, la sesión se aborta con error terminal.
    \item \textbf{Verificación en Lista Blanca (\textit{Whitelist}):} El orquestador consulta al \textit{Tenant B} para verificar si el \textit{Tenant A} figura en su lista de confianza permitida. Si la verificación resulta desfavorable, la comunicación se clausura impidiendo accesos no autorizados.
    \item \textbf{Entrega del Catálogo de Agent Cards:} Una vez aprobada la conexión, el \textit{Tenant B} suministra su catálogo de tarjetas de agente (\textit{Agent Cards}) disponibles. El orquestador transmite dicho catálogo al \textit{Agente A}.
    \item \textbf{Selección de Habilidad y Validación Cruzada:} El \textit{Agente A} inspecciona el catálogo y selecciona la tarjeta de agente idónea para la transacción. Si el agente requerido no está disponible, el flujo finaliza con estado \texttt{DECLINE}. Si existe, se emite una orden de selección (\texttt{SELECT}), tras la cual ambas partes ejecutan una verificación mutua de credenciales con el \textit{Agente B} antes de dar por completado el apretón de manos y transicionar al estado operativo inicial \texttt{on\_board}.
\end{enumerate}

\begin{figure}[htbp]
    \centering
    \includegraphics[width=0.85\textwidth,height=0.72\textheight,keepaspectratio]{handshake_sequence.png}
    \caption{Secuencia de Handshake bilateral: autenticación de tenants, validación en lista blanca (\textit{whitelist}), distribución del catálogo de Agent Cards y verificación cruzada de credenciales criptográficas.}
    \label{fig:handshake_flow}
\end{figure}

\subsection{Métodos JSON-RPC 2.0 Expuestos por el Gateway}
El Gateway implementa cuatro primitivas principales de negociación:
\begin{itemize}
    \item \texttt{create\_task}: Inicializa una sesión de negociación multilateral o bilateral. Recibe el identificador del tenant, la habilidad solicitada y el payload inicial.
    \item \texttt{send\_message}: Transmite un turno de negociación conteniendo el \texttt{facts\_delta} cuantitativo y el mensaje contextual hacia la contraparte.
    \item \texttt{get\_task}: Consulta el estado actual de la tarea, el emisor del turno activo y el resumen factual consolidado en Redis.
    \item \texttt{cancel\_task}: Provoca la cancelación inmediata de una negociación activa, transicionando a \texttt{frozen} (\texttt{failed}).
\end{itemize}

\subsection{Pipeline de Validación en Cascada en 4 Fases}
A fin de asegurar la robustez del sistema frente a mensajes malformados o ataques perimetrales, el Gateway procesa cada petición a través de cuatro filtros secuenciales:
\begin{enumerate}
    \item \textbf{Fase 1: Validación de Transporte y Protocolo Base:} Comprueba método HTTP POST, cabeceras requeridas (\texttt{application/json}) y conformidad estructural JSON-RPC 2.0.
    \item \textbf{Fase 2: Conformidad de Esquema A2A:} Valida que los parámetros correspondan a un \texttt{TaskRequest} o \texttt{MessagePart} conforme al metamodelo A2A.
    \item \textbf{Fase 3: Validación Semántica de Agent Card:} Verifica que la habilidad solicitada exista en la Agent Card activa y satisfaga su esquema de entrada.
    \item \textbf{Fase 4: Coherencia de Ciclo de Vida y Turnos:} Consulta atómicamente en Redis el estado de la tarea y la titularidad del turno activo para evitar operaciones desincronizadas.
\end{enumerate}

\subsection{Seguridad Perimetral y Control de Acceso}
El Gateway aplica directivas perimetrales antes de procesar el payload:
\begin{itemize}
    \item \textbf{Autenticación mTLS:} Validación mutua de certificados X.509 para comprobar la identidad del tenant emisor.
    \item \textbf{Rate Limiting y Cuotas:} Implementación de limitación de tasa por IP y por Tenant ID mediante algoritmo Token Bucket en Redis, previniendo denegación de servicio.
\end{itemize}
